\begin{frame}{13.1 자주 하는 실수}
  \begin{enumerate}
    \item \textbf{``시뮬레이션에서 검증했다 = 완료''로 선언}.
      시뮬레이션 점수는 \emph{시뮬레이션 환경}의 점수일 뿐 --- 시뮬레이션
      \emph{하위권} 정책이 실물에서 상위권으로 반전될 수 있음. 배포 전
      \emph{실물}에서(소량이라도) 평가하는 단계가 \textbf{반드시} 필요.
    \item \textbf{``시뮬레이션은 무한히 빠르니까 데이터 부족은 없는
      문제''}. 학습은 사실상 무한 데이터지만 \emph{배포}는 실물에서.
      ``시뮬레이션 1천만 에피소드''와 ``실물 0시간''은 \emph{다른 자원}.
    \item \textbf{무작위화 범위를 ``수학적으로 안전한'' 극단까지 확장}.
      $\mu = 100$ 예처럼 \emph{실물에서 불가능한} 조건까지 포함하면 정책이
      그 극단에 최적화되어 실물에서 오히려 나빠짐. ``넓을수록 안전''이
      \emph{아님}.
    \item \textbf{PD의 ``마법 수''를 RL로 대체할 수 있다고 생각}.
      $k_p > 5$ 는 이 진자의 \emph{물리}에서 유도되는 한계. RL이 이 한계를
      \emph{위배}하는 방향으로 학습하면 \textbf{학습 설정이 잘못됐다}는
      진단 신호. 물리 직관은 RL의 대체재가 아니라 \textbf{감시자}.
  \end{enumerate}
\end{frame}
